\section{Feature-Space Impact}
\label{sec:feature_space_impact}

Feature-space impact measures how much a transition changes the learned latent representation.
For transition $(o_t,a_t,o_{t+1})$, the impact term is defined as
\begin{equation}
    I_t=\left\lVert z_{t+1}-z_t\right\rVert_2.
    \label{eq:feature_space_impact}
\end{equation}
A transition with small $I_t$ leaves the representation nearly unchanged, while a transition with large $I_t$ moves the agent to a different region of latent space.
This quantity is related to impact-driven exploration, where exploration is encouraged through meaningful changes in learned features rather than only through novelty counts \cite{raileanu-2020}.

Impact is not equivalent to reward, novelty, or prediction error.
A transition can be novel but have little behavioral consequence, such as a small perturbation in an irrelevant observation component.
A transition can also have high prediction error because of stochasticity while not producing a useful change in controllable state.
The role of $I_t$ is to favor actions that induce represented state change, thereby complementing the learning-progress signal defined from prediction-error reduction.

The use of a learned feature space is essential to the definition.
If impact were computed directly in raw observation space, visually or numerically large differences could dominate the signal even when they are not relevant to control.
Computing impact in $\mathbb{R}^{d}$ instead makes the signal depend on the representation learned by the dynamics module.
The inverse-model component of the representation objective further encourages the feature space to retain information about action-dependent transitions.

Because the scale of $I_t$ can vary across environments and throughout training, impact is treated as a component of a normalized intrinsic score rather than as a standalone reward.
Let $v_I$ be an exponential moving estimate of the second moment of $I_t$,
\begin{equation}
    v_I \leftarrow \beta_{\mathrm{rms}}v_I+(1-\beta_{\mathrm{rms}})I_t^2,
\end{equation}
with $0<\beta_{\mathrm{rms}}<1$.
The normalized impact component is
\begin{equation}
    \widetilde{I}_t=\frac{I_t}{\sqrt{v_I+\varepsilon}},
\end{equation}
where $\varepsilon>0$ prevents division by zero.
This normalization makes the impact term comparable to the learning-progress term in the combined intrinsic reward.